\newpage
\setcounter{page}{2}
\renewcommand{\contentsname}{СОДЕРЖАНИЕ}
\tableofcontents

\newpage
\centeredheading{ВВЕДЕНИЕ}
\addcontentsline{toc}{chapter}{ВВЕДЕНИЕ}

Сегодня современные центры мониторинга информационной безопасности (SOC) сталкиваются с критической проблемой --- нехваткой квалифицированных специалистов и экспоненциальным ростом количества ложных срабатываний.
Для решения этой проблемы мировая индустрия ИБ двигается в сторону автоматизации реагирования на инциденты (SOAR) с применением больших языковых моделей — LLM\@.
Создание автономных ИИ-агентов, способных самостоятельно классифицировать инциденты, присваивать уровень угрозы, обогащать контекстом и применять сценарии сдерживания, является одним из самых актуальных R\&D направлений в кибербезопасности.

\textbf{Цель учебной практики:} Участие в исследовании, разработке и тестировании прототипа автоматизированной системы реагирования на инциденты информационной безопасности.
Проект «AutoSOAR».

В соответствии с индивидуальным заданием, практика была разделена на два основных этапа:
\begin{enumerate}
  \item Подготовка рабочей области и инфраструктуры проекта.
  \item Решение прикладных задач: адаптация модуля LLM-классификатора и разработка сценариев кибератак для тестирования системы на киберполигоне.
\end{enumerate}